% ============================================================================
% FILE: sections/03_implementation.tex
% PURPOSE: การติดตั้งจริง Mechanical / Electronics / STM32 Firmware
% ============================================================================

\section{การดำเนินการสร้างและติดตั้งระบบ (Implementation)}

% ============================================================
\subsection{การประกอบเชิงกล (Mechanical Assembly)}
% ============================================================

\subsubsection{ลำดับขั้นตอนการประกอบ}
% อธิบายลำดับการประกอบจากฐานถึง End Effector

\subsubsection{ผลการประกอบชิ้นส่วน}

% \begin{figure}[H]
%   \centering
%   \includegraphics[width=0.75\textwidth]{images/assembly_photo.jpg}
%   % [ใส่ภาพถ่ายจริงของ Assembly ที่ประกอบเสร็จ --- ด้านหน้า/ด้านข้าง]
%   \caption{หุ่นยนต์หลังการประกอบเสร็จสมบูรณ์}
%   \label{fig:assembly_photo}
% \end{figure}

\subsubsection{การตรวจสอบขนาดจริง}
% ผลการวัดขนาดฐานจริง เทียบกับ M1 (<= 200x200 mm)

\begin{table}[H]
  \centering
  \caption{ผลการตรวจสอบขนาดชิ้นส่วนจริง}
  \label{tab:mech_dim_check}
  % [ตาราง: รายการ | ค่าออกแบบ | ค่าวัดจริง | Req ID | Pass/Fail]
\end{table}

% ============================================================
\subsection{การประกอบตู้ควบคุมไฟฟ้า (Electrical Cabinet Assembly)}
% ============================================================

\subsubsection{รายการอุปกรณ์ภายในตู้}

\begin{table}[H]
  \centering
  \caption{Bill of Materials (BOM) ของตู้ควบคุม}
  \label{tab:elec_bom}
  % [ตาราง: อุปกรณ์ | รุ่น | จำนวน | หน้าที่]
\end{table}

\subsubsection{ผลการประกอบตู้ควบคุม}
% สถานะ: กำลังดำเนินการ --- จะอัปเดตหลังประกอบเสร็จ

% \begin{figure}[H]
%   \centering
%   \includegraphics[width=0.70\textwidth]{images/cabinet_interior.jpg}
%   % [ใส่ภาพถ่ายภายในตู้ควบคุมหลังประกอบเสร็จ]
%   \caption{ภายในตู้ควบคุมหลังการประกอบ}
%   \label{fig:cabinet_interior}
% \end{figure}

% \begin{figure}[H]
%   \centering
%   \includegraphics[width=0.70\textwidth]{images/cabinet_front.jpg}
%   % [ใส่ภาพถ่ายหน้าตู้ แสดง Power Indicator, E-Stop, Mode Switch]
%   \caption{หน้าตู้ควบคุมแสดงสวิตช์และไฟแสดงสถานะ}
%   \label{fig:cabinet_front}
% \end{figure}

% ============================================================
\subsection{สถาปัตยกรรม STM32 Firmware (Firmware Architecture)}
% ============================================================

\subsubsection{ภาพรวมสถาปัตยกรรม Firmware}
% อธิบายว่า Firmware ทำงานอย่างไรโดยรวม
% ระบุ Main Loop, Interrupt, Task ต่างๆ

% \begin{figure}[H]
%   \centering
%   \includegraphics[width=0.85\textwidth]{images/firmware_arch_block.png}
%   % [ใส่ Block Diagram ของ Firmware: Main Loop, ISR, Module ต่างๆ]
%   \caption{ภาพรวมสถาปัตยกรรม STM32 Firmware}
%   \label{fig:fw_arch}
% \end{figure}

\subsubsection{โครงสร้างไฟล์และโมดูล (Code Structure)}
% Folder Structure และหน้าที่ของแต่ละ Module

\begin{table}[H]
  \centering
  \caption{โมดูลหลักใน STM32 Firmware}
  \label{tab:fw_modules}
  % [ตาราง: ชื่อไฟล์ | หน้าที่ | Peripheral ที่ใช้ | สถานะ]
\end{table}

\subsubsection{Interrupt Service Routine (ISR) และ Timing}
% ISR ที่ใช้ เช่น Encoder ISR, Timer ISR สำหรับ PID Loop

\begin{table}[H]
  \centering
  \caption{ISR และ Task Timing ของ Firmware}
  \label{tab:fw_timing}
  % [ตาราง: ISR/Task | Trigger | Period/Priority | หน้าที่]
\end{table}

\subsubsection{การทำงานร่วมกันของโมดูล (Module Interaction)}
% Encoder -> PID -> PWM Output

% \begin{figure}[H]
%   \centering
%   \includegraphics[width=0.85\textwidth]{images/firmware_dataflow.png}
%   % [ใส่ Data Flow Diagram ระหว่าง Module ต่างๆ ใน Firmware]
%   \caption{Data Flow ระหว่างโมดูลใน Firmware}
%   \label{fig:fw_dataflow}
% \end{figure}

\subsubsection{การ Implement S-Curve Motion Profile ใน Firmware}
% วิธีการแปลง S-Curve ที่คำนวณได้มาเป็น Code จริงบน STM32
% อธิบาย State Machine ของ Motion Profile

\subsubsection{Joystick Mode และการ Map ปุ่ม}
% โหมดที่ใช้งานได้ทั้งหมด: Manual, Auto, Homing
% อ้างอิง J1, E8

\begin{table}[H]
  \centering
  \caption{การ Map ปุ่ม Joystick กับคำสั่งระบบ}
  \label{tab:joystick_map}
  % [ตาราง: ปุ่ม | โหมด | คำสั่งที่ส่ง | Req ID]
\end{table}